\subsection{准希尔伯特空间}

定义 4. --- 准希尔伯特空间是一个赋有域 K 上的向量空间结构和一个正埃尔米特型的集合 E。当 K 是 R（相应地，K 是 C）时，我们说 E 是实（相应地，复）准希尔伯特空间。

例. --- 1) 型 \((\lambda, \mu) \mapsto \overline{\lambda}\mu\) 在 K 上定义了一个准希尔伯特结构，称为典范结构。当把 K 看作准希尔伯特空间时，除非另有说明，我们总指它具有这个结构。

\begin{enumerate}
\def\labelenumi{\arabic{enumi})}
\setcounter{enumi}{1}
\tightlist
\item
  设 I 为 R 中的一个区间（有界或无界），E 为定义在 I 上、取值于 C 且具有紧支撑的受控函数（FVR, II, p.~4）组成的集。显然
\end{enumerate}

\(\mathbf{E}\) 是 \(\mathbf{C}\) 上的向量空间；设 \(f\) 为半双线性型 \((x, y) \mapsto \int_{1}^{\infty} \overline{x(t)} y(t) dt\) ；

立即可以看出 \(f\) 是 \(E\) 上的正埃尔米特型，从而在此空间上定义了一个准希尔伯特结构。

\begin{enumerate}
\def\labelenumi{\arabic{enumi})}
\setcounter{enumi}{2}
\tightlist
\item
  设 \(n \geqslant 0\) 为整数。我们在空间 \(\mathbf{K}^n\) 上借助埃尔米特型
\end{enumerate}

\[
(x, y) \mapsto \sum_ {j = 1} ^ {n} \overline {{{{x}}}} _ {j} y _ {j}
\]

定义一个准希尔伯特空间结构（对 \(x = (x_{1}, \ldots, x_{n})\) 与 \(y = (y_{1}, \ldots, y_{n})\) ）。当 \(\mathbf{K}\) 是 \(\mathbf{R}\) 时，我们看到这正是 \(\mathbf{R}^{n}\) 中两个向量的内积（GT, VI, § 2, No.~2）。

* 4) 设 \(\ell^2\) （或 \(\ell^2(\mathbf{N})\) ）为序列 \(x = (x_n)_{n \in \mathbb{N}}\) （元素在 \(\mathbf{K}\) 中且使 \(\sum_{n=0}^{\infty} |x_n|^2\) 为有限）组成的集。可以证明 \(\ell^2\) 是 \(\mathbf{K}^{\mathbf{N}}\) 的向量子空间，并可在其上定义一个准

希尔伯特空间结构于 \(\ell^2\) ，借助埃尔米特型 \((x, y) \mapsto \sum_{n=0}^{\infty} \overline{x}_n y_n\) （参见~V, p.~18）。

\begin{enumerate}
\def\labelenumi{\arabic{enumi})}
\setcounter{enumi}{4}
\tightlist
\item
  设 \(\mathbf{E}\) 为实准希尔伯特空间，\(f\) 为 \(\mathbf{E}\) 上相应的对称双线性型。设 \(\mathrm{E}_{(\mathbf{C})}\) 为 \(\mathbf{E}\) 的向量空间复化；把 \(\mathbf{E}\) 等同于 \(\mathrm{E}_{(\mathbf{C})}\) 的一个子集（借助映射 \(x \mapsto 1 \otimes x\) ），使 \(\mathrm{E}_{(\mathbf{C})}\) 的每个元素都可唯一地写成 \(x_1 + ix_2\) ，其中 \(x_1, x_2\) 属于 \(\mathbf{E}\) 。映射 \(f\) 唯一地延拓为一个埃尔米特型 \(f_{(\mathbf{C})}\) （\(\mathrm{E}_{(\mathbf{C})}\) 上的）；我们有，
\end{enumerate}

\[
f _ {(\mathbf {C})} \left(x _ {1} + i x _ {2}, y _ {1} + i y _ {2}\right) = f \left(x _ {1}, y _ {1}\right) + f \left(x _ {2}, y _ {2}\right) + i \left(f \left(x _ {1}, y _ {2}\right) - f \left(x _ {2}, y _ {1}\right)\right).
\]

特别地，我们有

\[
f _ {(\mathbf {C})} (x _ {1} + i x _ {2}, x _ {1} + i x _ {2}) = f (x _ {1}, x _ {1}) + f (x _ {2}, x _ {2}) \geqslant 0,
\]

因此 \(f_{(\mathbf{C})}\) 是正的。我们说 \(\mathrm{E}_{(\mathbf{C})}\) （赋有 \(f_{(\mathbf{C})}\) ）是 \(\mathbf{E}\) 的准希尔伯特空间复化。

每当只考虑向量空间 E 上的一个准希尔伯特空间结构时，对 E 的一对点 \((x, y)\) ，定义该结构的埃尔米特型在这对点上的值记作 \(\langle x|y\rangle_{E}\) ，或在不致混淆时简记为 \(\langle x|y\rangle\) 。这个数称为 x 与 y 的内积 \(^{1}\) （当 y = x 时称为 x 的标量平方）。向量 x, y 称为正交的，如果 \(\langle x|y\rangle = 0\) 。函数 \(x \mapsto \|x\| = \langle x|x\rangle^{1/2}\) 是向量空间 E 上的半范数（V, p.~3）；准希尔伯特空间总是被认为赋有这个半范数（从而也赋有相应的拓扑与一致结构）。

采用这些记号，在准希尔伯特空间 E 中，Cauchy-Schwarz 不等式可写成

\[
\left| \langle x | y \rangle \right| \leqslant \| x \|\cdot \| y \|.\tag{9}
\]

因此，内积是 \(\mathbf{E} \times \mathbf{E}\) 上的连续半双线性型（II, p.~5, 命题 4）。

为使 E 是 Hausdorff 的，必要且充分的是 \(x \mapsto \|x\|\) 是 E 上的范数；换言之，埃尔米特型 \((x, y) \mapsto \langle x | y \rangle\) 是正的且非退化的；这等价于说 0 是 E 中唯一与自身正交的向量。

按照一般定义（S, IV, § 1, No.~5），准希尔伯特空间 E 到准希尔伯特空间 F 上的同构是从 E 到 F 上的一个双射线性映射 u，它满足

\[
\langle u (x) | u (y) \rangle = \langle x | y \rangle\tag{10}
\]

对 E 中每个 x 与 y。由此推出，\(\|u(x)\| = \|x\|\) 对所有 \(x \in E\) 成立，且 u 显然是 E 与 F 的拓扑向量空间结构之间的同构；若 E 与 F 是 Hausdorff 的，则 u 是 E 到 F 上的一个等距。反之，若 u 是从 E 到 F 上的双射线性映射，且 \(\|u(x)\| = \|x\|\) 对所有 \(x \in E\) 成立，则极化公式（V, p.~2）表明 u 是 E 到 F 上的准希尔伯特空间同构。

设 \(E\) 为复准希尔伯特空间，\(\langle x|y\rangle\) 为 \(E\) 中的内积。在集合 \(E\) 上，可以定义关于 \(C\) 的第二个向量空间结构，取同样的加法群规律，并以 \((\lambda, x) \mapsto \overline{\lambda}x\) 为外合成规律（A, II, § 1, No.~13）；对这个向量空间结构，\((x, y) \mapsto \langle y|x\rangle\) 是一个正埃尔米特

\(^{1}\) 有时我们会用 \((x|y)\) 表示 \(\langle y|x\rangle\) 。注意 V, p.~2 的公式 (4) 取下列等价形式：

(4')

\[
\langle \sum_ {i} \lambda_ {i} x _ {i} | \sum_ {j} \mu_ {j} y _ {j} \rangle = \sum_ {i, j} \overline {{{{\lambda}}}} _ {i} \mu_ {j} \langle x _ {i} | y _ {j} \rangle .\tag{4"}
\]

\[
(\sum_ {i} \lambda_ {i} x _ {i} | \sum_ {j} \mu_ {j} y _ {j}) = \sum_ {i, j} \lambda_ {i} \overline {{\mu}} _ {j} (x _ {i} | y _ {j}).
\]

型。给 E 赋以这个新的向量空间结构与新的埃尔米特型所得的准希尔伯特空间 \(\overline{E}\) 称为 E 的共轭空间。E 到 \(\overline{E}\) 上的同构 u 是 E 到自身的一个半线性映射（关于 C 的自同构 \(\xi \mapsto \overline{\xi}\) ），满足 \(\langle u(y)|u(x)\rangle = \langle x|y\rangle\) 或 \(\langle u(x)|u(y)\rangle = \overline{\langle x|y\rangle}\) （对 E 中的 x, y）；这样的映射称为准希尔伯特空间 E 的半自同构。

若 E 是准希尔伯特空间，M 是 E 的向量子空间，则内积 \(\langle x|y\rangle\) 在 \(M \times M\) 上的限制是 M 上的正埃尔米特型，从而在 M 上定义了一个准希尔伯特空间结构；我们说这个结构由 E 的结构诱导，或说 M 是 E 的准希尔伯特子空间。

